% ===========================================================
% Conference Paper Starter Template
% ===========================================================
% Replace the content below with your own title, authors,
% abstract, sections, and citations.
% ===========================================================

\documentclass[11pt]{article}

% ---- Packages ----
\usepackage[T1]{fontenc}
\usepackage[utf8]{inputenc}
\usepackage[a4paper,margin=1in]{geometry}
\usepackage{natbib}
\usepackage{amsmath,amssymb,amsfonts}
\usepackage{graphicx}
\usepackage{textcomp}
\usepackage{xcolor}
\usepackage{hyperref}
\usepackage{setspace}

\graphicspath{{figures/}}
\onehalfspacing

% ===========================================================
\begin{document}

\title{The Corporate Stage: Algorithmic Surveillance, Hybrid Visibility, \\
and the Paradox of Spectatorship in the Modern Workplace}

\author{
  Vivien Jiaqian Zhu \\
  \textit{Department of East Asian Languages and Cultures, University of California, Berkeley} \\
  \textit{Visiting Scholar, Stanford University} \\
  \textit{Research Fellow, University of Cambridge} \\
  \texttt{email@example.com}
}

\date{}

\maketitle

% ===========================================================
\begin{abstract}
% TODO: 150--250 word abstract for IntechOpen submission.
% Should state: (1) the problem -- the convergence of algorithmic
% monitoring and performative self-presentation in contemporary
% hybrid/remote workplaces; (2) the theoretical apparatus -- Goffman's
% dramaturgical model of everyday performance, Foucault's account of
% disciplinary and panoptic power, and Zuboff's theory of surveillance
% capitalism; (3) the central claim -- that hybrid visibility produces
% a "paradox of spectatorship," in which workers are simultaneously
% audience and performer, watched and self-watching, under conditions
% that exceed the classical panopticon; (4) what the chapter
% contributes to workplace studies / organizational theory.
\end{abstract}

\textbf{Keywords:} algorithmic surveillance, hybrid work, dramaturgical
theory, panopticism, surveillance capitalism, spectatorship, workplace
visibility

% ===========================================================
\section{Introduction}
\label{sec:intro}

% TODO: Open with the concrete phenomenon -- webcam monitoring software,
% keystroke logging, "productivity scores," always-on video calls -- as
% the empirical entry point into the theoretical argument. Motivate why
% the shift to hybrid/remote work makes visibility a live analytic
% problem rather than a background condition of employment.

\subsection{The Corporate Stage as Analytic Frame}
% TODO: Introduce the chapter's central metaphor -- the "corporate
% stage" -- as a deliberate extension of Goffman's dramaturgy into a
% context he did not theorize: algorithmically mediated, asymmetric,
% and persistent observation.

\subsection{Contributions / Chapter Roadmap}
The chapter makes the following moves:
\begin{itemize}
  \item Synthesizes Goffman, Foucault, and Zuboff into a single
        framework for reading hybrid workplace visibility.
  \item Names and develops the ``paradox of spectatorship'' as a
        distinct condition of contemporary algorithmic management.
  \item % TODO: third contribution -- e.g., implications for worker
        % autonomy, organizational design, or policy.
\end{itemize}

% ===========================================================
\section{Theoretical Framework}
\label{sec:framework}

\subsection{Goffman: The Presentation of Self as Front-Stage Labor}
% TODO: Goffman's front-stage/back-stage distinction, impression
% management, and what happens when the back-stage collapses --
% i.e., when home becomes the site of a permanently front-stage
% performance under webcam surveillance. Cite~\citep{goffman1959}.

\subsection{Foucault: Panopticism and the Internalization of the Gaze}
% TODO: Foucault's panopticon as a disciplinary architecture that
% works through the possibility, not certainty, of being watched --
% and how algorithmic monitoring both intensifies and departs from
% this model (continuous data capture vs. the panopticon's
% uncertainty principle). Cite~\citep{foucault1975}.

\subsection{Zuboff: Surveillance Capitalism and Behavioral Data Extraction}
% TODO: Zuboff's account of behavioral surplus and the commodification
% of experience, extended from consumer contexts into the workplace as
% a site of data extraction from labor itself. Cite~\citep{zuboff2019}.

\subsection{Toward a Synthesis: Hybrid Visibility}
% TODO: Note also relevant empirical/educational-technology literature
% on dashboards, learning analytics, and visibility as a design
% problem -- e.g.,~\citep{rets_TODO}. [Need full citation from Vivien --
% could not verify exact title/venue via search.]

% ===========================================================
\section{Hybrid Visibility in the Modern Workplace}
\label{sec:hybrid}

\subsection{From Panopticon to Platform: What Changes}
% TODO: Distinguish classical disciplinary visibility (architectural,
% spatial, uncertain) from platform-based visibility (continuous,
% quantified, data-retentive, algorithmically scored).

\subsection{The Home as Front Stage}
% TODO: The domestic space reconfigured as a site of professional
% performance; blurred back-stage/front-stage boundaries under
% webcam-based and activity-tracking monitoring regimes.

\subsection{Metrics, Scores, and the Datafied Self}
% TODO: Productivity dashboards, "active time," algorithmic scoring
% as a new register of visibility distinct from direct human
% observation.

% ===========================================================
\section{The Paradox of Spectatorship}
\label{sec:paradox}

\subsection{Audience and Performer at Once}
% TODO: Core theoretical move of the chapter -- the worker as
% simultaneously spectator of their own datafied self (via dashboards,
% self-monitoring apps) and performer for an algorithmic/managerial
% audience. Unpack the reflexive loop this creates.

\subsection{Illustrative Cases}
% TODO: 2--3 concrete illustrative cases or vignettes (e.g., call
% center monitoring, remote-work productivity software, gig-platform
% rating systems) to ground the theoretical argument. Confirm whether
% IntechOpen chapter format expects a dedicated "case studies" section
% or integrated examples.

\subsection{Limits of the Paradox: Resistance and Opacity}
% TODO: Worker tactics of partial opacity, gaming of metrics, or
% refusal -- as a counterpoint to a purely totalizing surveillance
% narrative.

% ===========================================================
\section{Discussion}
\label{sec:discussion}

% TODO: Implications for organizational design, management practice,
% and worker wellbeing. Position the chapter's contribution within the
% edited volume's broader theme (\textit{The Evolving Landscape of
% Workplace -- Challenges and Opportunities}, ed. C\u{a}t\u{a}lina Radu).

% ===========================================================
\section{Conclusion}
\label{sec:conclusion}

% TODO: Restate the paradox of spectatorship as the chapter's central
% contribution; gesture toward future research (e.g., comparative
% regulatory approaches, worker-side sousveillance tools).

% ===========================================================
\bibliographystyle{apalike}
\bibliography{references}

\end{document}
